\section*{Chapitre \ref{ch:pl:real-time-risk} --- Risque en temps réel}

\begin{solution}{exo:pl:real-time-risk:1}
$-0.4 \times 1\,000 \times 5\,000 = -\$2$~millions de delta. Sous 90~\% de 50~millions de dollars (45~millions de dollars), 22 transactions de ce type tiennent encore, avant que le gamma et le
mouvement du marché ne modifient les deltas.
\end{solution}

\begin{solution}{exo:pl:real-time-risk:2}
La bibliothèque calcule le delta et le gamma par différences centrées, avec des perturbations de $\pm1~\%$ du spot, si bien que l'approximation quadratique $\Delta\,\mathrm{d}S +
\frac12\Gamma\,\mathrm{d}S^2$ passe exactement par les deux points réévalués~: à $\pm1~\%$, elle est la réévaluation complète.
\end{solution}

\begin{solution}{exo:pl:real-time-risk:3}
Que son desk resterait à $-\$45.2$~millions, sous sa limite, et que la firme passerait à $-\$92.8$~millions, en dépassant celle de la firme. Le système du desk~B
ne connaît que les transactions du desk~B~; le chiffre de la firme exige aussi celles du desk~A.
\end{solution}

\begin{solution}{exo:pl:real-time-risk:4}
L'approximation quadratique garde le gamma constant. Le gamma d'une option de vente achetée baisse quand l'indice monte et que l'option passe hors de la monnaie, et sa valeur s'aplatit
vers zéro, si bien que l'approximation quadratique, qui reste courbe, fausse de plus en plus la perte de valeur à la hausse~; à la baisse, les puts entrent dans la
monnaie et se comportent de façon plus linéaire, plus proche de l'approximation quadratique.
\end{solution}

\begin{solution}{exo:pl:real-time-risk:5}
Acceptable tant que le spot reste dans le seuil de rafraîchissement du spot du cache (le delta reporté est alors exact à
l'erreur du seuil près) et que rien d'autre dont dépendent les sensibilités n'a bougé~; inacceptable après un choc de volatilité, un fort
mouvement d'un facteur corrélé, un changement du livre que le cache n'a pas vu, ou le temps qui passe près de l'échéance.
\end{solution}

\begin{solution}{exo:pl:real-time-risk:6}
Parce que l'agrégateur ne recalcule jamais l'arbre, une erreur dans un écart --- une contribution non mise à jour, une ancienne valeur soustraite deux fois,
un ancêtre omis --- persiste et s'accumule en silence. Comparer avec une reconstruction après exécutions, ticks et rafraîchissements permet de détecter toute dérive.
\end{solution}

\begin{solution}{exo:pl:real-time-risk:7}
À 0.1~\% le cache est réévalué 105 fois, à 2~\% une seule fois~; l'exposition de la firme à la clôture est de $-\$90.837$~millions et $-\$90.832$~millions, et le
moment du dépassement est le même, 14:23:50. Pour ce livre et cet après-midi, le seuil modifie le coût, non la réponse~; un mouvement plus grand ou un
livre à échéance plus proche changerait cela.
\end{solution}

\begin{solution}{exo:pl:real-time-risk:8}
Les limites de desk bornent chaque desk au moment de chaque transaction~; les expositions dérivent avec le marché après les transactions (à cause du gamma), et les limites sont souvent fixées de sorte que
leur somme dépasse l'appétit de la firme, et même lorsqu'elle reste en deçà, les expositions des desks sur des sous-jacents corrélés peuvent s'accumuler alors que
la compensation dans l'arithmétique des limites ne le fait pas. Seul un agrégat calculé sur l'ensemble des desks, en continu, donne le chiffre de la firme.
\end{solution}

\begin{solution}{pb:pl:real-time-risk:1}
\begin{enumerate}
\item Des secondes et la firme entière, entre le moteur nocturne (heures, réévaluation complète) et le contrôle avant envoi (microsecondes, un ordre).
\item Les delta-dollars des transactions sont additionnés dans la hiérarchie de risque~; on ne s'attend pas à ce que les desks utilisent leurs limites ensemble, si bien que l'appétit de la firme est inférieur
à leur somme.
\item Une exécution ajoute à chaque ancêtre la différence entre sa nouvelle contribution et son ancienne~; un tick ajoute, par transaction, la variation de ses delta-dollars reportés.
\item Lorsque le spot a bougé de plus de 0.5~\% depuis le remplissage du cache~; en réévaluant chaque instrument avec la bibliothèque de valorisation.
\item Les expositions, l'heure de marché, le spot et l'âge du cache, et les limites dépassées.
\item Le desk~A a acheté des puts, 1\,000 à la fois~; le desk~B a vendu des contrats à terme, 400 à la fois~; chacun a vérifié que sa propre exposition restait sous 90~\% de sa limite.
\item $-\$49.3$~millions et $-\$43.6$~millions.
\item À 14:23:50.
\item $-\$90.8$~millions, soit 1.82 fois l'appétit de la firme.
\item Desk~B sous la limite à $-\$45.2$~millions, la firme à $-\$92.8$~millions, en dépassement.
\item P\&L delta-gamma issu des grecques calculées à la clôture, comparé à la variation des valeurs complètes, pour des mouvements de $-10~\%$ à $+10~\%$.
\item $-6.75~\%$ et $+3.5~\%$.
\item 3.1~\% et 23.6~\%.
\item Les grecques ont été calculées à partir de perturbations de $\pm1~\%$.
\item Elles proviennent de la réévaluation complète.
\item \textbf{Résultat principal.} La limite de la firme a été dépassée à 14:23:50, alors qu'aucun des deux desks n'a dépassé la sienne (les valeurs extrêmes de $-\$49.3$~millions et
$-\$43.6$~millions, par rapport à 50~millions de dollars)~; à la clôture la firme était à $-\$90.8$~millions, soit 1.82 fois son appétit~; l'erreur de P\&L du cache dépasse
1~\% de la réévaluation complète au-delà d'un mouvement de $+3.5~\%$ (ou $-6.75~\%$).
\item Un risk manager doté de l'autorité sur les deux desks, capable de bloquer au contrôle avant envoi les ordres qui accroissent le risque et de réduire une limite via le
\omterm{def:pl:signal-serving-and-parameter-management:store}{magasin de paramètres}.
\item À partir de l'erreur d'approximation~: assez petite pour que l'erreur reste nettement sous la précision retenue pour les limites (ici, sous 1~\% jusqu'à 3.5~\%),
et assez grande pour que les coûts de réévaluation restent bornés.
\item Vega, sensibilités aux taux et au crédit, notionnels bruts et nets, P\&L comparé aux limites de perte, et concentrations par émetteur.
\item Parce qu'une limite franchie à l'échelle des desks se forme pendant qu'ils négocient, et n'est utile à connaître que tant qu'on peut encore l'arrêter.
\end{enumerate}
\end{solution}

\begin{solution}{iq:pl:real-time-risk:1}
Les limites de desk sont fixées séparément pour chaque desk et s'additionnent en général à plus que l'appétit de la firme~; les expositions de même sens et de même
sous-jacent s'additionnent, et le gamma fait bouger les expositions après les transactions.
\iqlookfor{Agrégation entre les desks.}
\end{solution}

\begin{solution}{iq:pl:real-time-risk:2}
\omterm{def:pl:real-time-risk:incremental}{Agrégation incrémentale} dans une hiérarchie de risque~: chaque événement ajoute son écart à ses ancêtres~; sensibilités mises en cache et rafraîchies aux seuils~;
limites contrôlées sur les seuls nœuds mis à jour~; instantanés horodatés~; une reconstruction périodique comme contrôle.
\iqlookfor{Écarts, non recalcul~; caches avec invalidation.}
\end{solution}

\begin{solution}{iq:pl:real-time-risk:3}
Pour les mouvements dans la plage où son erreur reste petite au regard des décisions prises avec elle --- ici sous 1~\% jusqu'à 3.5~\% --- et pour les portefeuilles sans
de forts termes d'ordre supérieur~; réévaluer pour les grands mouvements, les scénarios de stress, les barrières et les options à courte échéance proches du strike.
\iqlookfor{Mesurer l'erreur en fonction de l'ampleur du mouvement.}
\end{solution}

\begin{solution}{iq:pl:real-time-risk:4}
L'heure de marché qu'elle reflète, l'heure et le point de marché des sensibilités utilisées et leur âge, quelles positions sont en attente, et quelles
limites sont dépassées.
\iqlookfor{Horodatages et fraîcheur des données.}
\end{solution}

\begin{solution}{iq:pl:real-time-risk:5}
Les exécutions du \omterm{def:pl:position-and-pnl-services:service}{service de positions} et les données de marché du flux alimentent un agrégateur sur la hiérarchie de risque~; les sensibilités issues de la bibliothèque de
valorisation sont placées dans un cache avec invalidation~; mises à jour incrémentales~; limites tirées d'un \omterm{def:pl:signal-serving-and-parameter-management:store}{magasin de paramètres}, contrôlées à chaque mise à jour, avec des alertes envoyées aux personnes qui peuvent
agir~; \omterm{def:pl:real-time-risk:whatif}{contrôles what-if} avant envoi des ordres~; instantanés avec leur fraîcheur~; réévaluation complète périodiquement et pour les stress.
\iqlookfor{Hiérarchie, incrémentalité, cache, limites dans la boucle.}
\end{solution}
